\section{Introduction}
\label{sec:intro}

A useful world model need not reproduce every pixel. It must preserve the distinctions that change
what an agent should do. For a robot approaching a doorway, the texture of the wall may be irrelevant;
the wall's position and the available opening are not. Learning this separation without specifying a
particular downstream reward is a central attraction of joint-embedding predictive architectures
(JEPAs): predict an informative representation of the future rather than reconstruct the future
image \citep{lecun2022path,assran2023ijepa,bardes2024vjepa}.

Recent systems make this approach increasingly practical. Pretrained video representations support
action-conditioned planning, and end-to-end objectives jointly learn representations and dynamics
\citep{assran2025vjepa2,sobal2025pldm,maes2026lewm}. These successes raise an evaluation question:
\emph{when a learned world model supports one behavior, what has that established about its usefulness
for another?} Nearby goal reaching, recognizable decoded rollouts, and reward-driven control measure
different capabilities. A representation can encode a feature that its predictor does not preserve;
a predictor can achieve low error on features that a controller cannot use.

We study these distinctions in a squared-loss JEPA with Gaussian marginal regularization, following
the end-to-end formulation of LeWorldModel \citep{maes2026lewm,balestriero2025lejepa}. The study moves
from a complex environment to progressively simpler ones (Figure~\ref{fig:staircase}). Craftax first
exposes a split between short-horizon goal reaching and reward control. A moving-wall task reduces
the decision to one observable obstacle and an initially hidden gap. Cube worlds then remove the
obstacle and vary only visual nuisance and dynamics. The simplification makes the question sharper:
is relevant information absent from the encoding, lost during prediction, or ineffective in the
complete control pipeline?

The gate provides the clearest empirical separation. A decoder recovers the wall from frozen real
encodings, and a convolutional policy trained on those decoded observations discovers the survival
strategy in two of nine runs. The same decoder applied to predicted latents loses the wall within
two steps for all three fitted encoders. None of the 27 policies trained in the tested JEPA
imagination pipeline discovers the strategy. The decoder adds supervision, and imagined training
also changes reward estimation and policy optimization; these distinctions are essential to what
the experiment establishes.

Our contributions are threefold. First, we organize the evidence around separate tests of
\textbf{retention, prediction, and control}, making clear which transfers between these tests are
supported. Second, we provide a controlled, replicated gate study showing useful information in a
frozen representation alongside poor performance of its imagination pipeline, with predictor
variants and model seeds reported separately. Third, we characterize two limits of the training
objective: Gaussian non-collapse can coexist with nuisance-only encoding, and exact decodability
can coexist with arbitrarily small prediction-loss scale. The predictable-nuisance mechanism is
known \citep{sobal2022slow}; our theory specializes it to Gaussian regularization and separates it
from the distinct question of how much loss is assigned to retained information.

\begin{figure}[t]
\centering
\paperfig{evidence_staircase.pdf}{1.0}
\caption{\textbf{Simplify the environment to sharpen the question.} Craftax tests whether local
visual-goal planning transfers to reward-driven behavior. The gate exposes a single control-relevant
obstacle and separates encoded observations from imagined futures. Cube worlds isolate predictable
nuisance. The environments supply complementary diagnostics, not a shared performance leaderboard.}
\label{fig:staircase}
\end{figure}
